\documentclass[11pt,a4paper]{article}
\usepackage[T1]{fontenc}
\usepackage{isabelle,isabellesym}
\usepackage{amssymb}
\usepackage{booktabs}

% this should be the last package used
\usepackage{pdfsetup}

% urls in roman style, theory text in math-similar italics
\urlstyle{rm}
\isabellestyle{it}

% ToC: widen the subsection number box (two-component numbers >= 10
% otherwise abut their titles) and the right margin reserved for page
% numbers (long generated headings otherwise collide with them).
\makeatletter
\renewcommand*\l@subsection{\@dottedtocline{2}{1.5em}{3.3em}}
\renewcommand\@tocrmarg{4.2em}
\makeatother

% Generated theory text carries long unbreakable identifiers; allow
% extra stretch before lines protrude into the margin.
\emergencystretch=3em

\begin{document}

\title{DBLog Virtual Cuts:\\
  Certified Snapshot-Equivalent Replay of Live Databases}
\author{Andreas Andreakis}
\maketitle

\begin{abstract}
  DBLog reconstructs a logical, replay-equivalent snapshot of a live
  database scope without taking a physical database snapshot, by
  interleaving chunked reads with change-data-capture (CDC) events.
  This entry machine-checks the theory behind that protocol, following
  the companion paper~\cite{andreakis2026theoretical}: it formalizes
  \emph{certified virtual cuts} --- finite, asynchronously gathered
  mixes of CDC events and chunk reads whose certified replay reaches
  the source state at a chosen frontier on a chosen key scope --- and
  proves the layered theorem ladder relating run wellformedness,
  verifier acceptance, and replay equality, together with a
  source-side continuation and restriction algebra and a whole-table
  specialization.

  The development is axiom-free: it contains no
  \isa{axiomatization}, \isa{typedecl}, or \isa{consts} declaration,
  and its only \isa{typedef} is a conservative copy of the naturals
  for source coordinates. All witnesses and fixtures are constructed
  models with locale interpretations, and a closing non-vacuity
  theory exercises every main theorem at concrete witnesses, the
  certificate-layer results at genuinely accepted certificates. Every theorem is conditional on the premises stated
  in its statement; the machine check does not discharge the
  deployment obligations or the external-observation assumption that
  a real deployment must establish.
\end{abstract}

\tableofcontents

\newpage

\section*{Entry overview}
\addcontentsline{toc}{section}{Entry overview}

\subsection*{What is proved}

The load-bearing spine is the Layer~2 result: the canonical clean
prefix of a wellformed DBLog run, replayed from the empty state,
reproduces the source state at the run's frontier restricted to the
run's scope (\isa{wellformed{\isacharunderscore}run{\isacharunderscore}implies{\isacharunderscore}virtual{\isacharunderscore}cut});
the base state enters through the chunk reads, not the replay seed.
On top of it sit an interface layer and an algebra:

\begin{itemize}
\item \emph{Layer 3 (certificates).} Inside the
  \isa{layer{\isadigit{3}}{\isacharunderscore}checker{\isacharunderscore}substrate}
  locale --- which abstracts the verifier as a set of soundness
  obligations --- an accepted certificate under faithful source
  observation replays to the source state at the certified frontier
  on the certified scope. The obligations are proved for the concrete
  verifier shipped with the entry, so the locale hypotheses are
  satisfiable.
\item \emph{Layer 4 (whole table).} The specialization of the Layer~3
  result to whole-table claim scopes: the replay reproduces the
  entire source state.
\item \emph{Continuation algebra.} Source-side continuation across
  frontiers (appending the faithful CDC segment for a half-open
  frontier interval), sub-scope restriction, the whole-table
  continuation instance, and an accepted-certificate continuation
  result.
\end{itemize}

By name, the nine main theorems are
\isa{wellformed{\isacharunderscore}run{\isacharunderscore}implies{\isacharunderscore}virtual{\isacharunderscore}cut}
(Layer~2);
\isa{accepted{\isacharunderscore}certificate{\isacharunderscore}implies{\isacharunderscore}wellformed{\isacharunderscore}run}
and \isa{accepted{\isacharunderscore}virtual{\isacharunderscore}cut{\isacharunderscore}sound}
(Layer~3, in the checker locale);
\isa{accepted{\isacharunderscore}whole{\isacharunderscore}table{\isacharunderscore}anchor{\isacharunderscore}domain{\isacharunderscore}specialization}
(Layer~4, in the checker locale); the four state-level
continuation/restriction results
\isa{virtual{\isacharunderscore}cut{\isacharunderscore}state{\isacharunderscore}continuation},
\isa{virtual{\isacharunderscore}cut{\isacharunderscore}state{\isacharunderscore}restrict{\isacharunderscore}scope},
\isa{whole{\isacharunderscore}table{\isacharunderscore}state{\isacharunderscore}continuation},
and
\isa{virtual{\isacharunderscore}cut{\isacharunderscore}restrict{\isacharunderscore}to{\isacharunderscore}subscope};
and
\isa{accepted{\isacharunderscore}certificate{\isacharunderscore}continuation{\isacharunderscore}sound}
(both in the theory \isa{Continuation}). The
closing theory fires each of the nine at concrete witnesses ---
the certificate-layer results at genuinely accepted certificates.

Composition is proved; single-coordinate chunk-read attribution and
CDC faithfulness are explicit premises, and the operational watermark
is not a formal object the theorems consult. The deployment
obligations and the external-observation assumption of the
paper~\cite{andreakis2026theoretical} appear here as explicit theorem
premises and locale obligations; the entry makes them visible rather
than discharging them. The protocol itself is from the original DBLog
system paper~\cite{andreakis2020dblog}. This document corresponds to
version~2.1 of the development, archived on
Zenodo~\cite{andreakis2026artifact}; its frozen predecessors remain
separately archived (v2.0 at DOI 10.5281/zenodo.20652511, v1.0 at DOI
10.5281/zenodo.20389697). The
development is distributed under the BSD 3-Clause license.

\subsection*{Trust base}

The session and its shared
\isa{Dual{\isacharunderscore}Write{\isacharunderscore}Layer{\isadigit{0}}}
parent contain no \isa{axiomatization}, \isa{typedecl}, or
\isa{consts} declaration. The development's single \isa{typedef} ---
in the parent theory \isa{Source{\isacharunderscore}History} ---
introduces the source-coordinate type as a copy of \isa{nat} (a
conservative extension over a provably nonempty set), with its
\isa{linorder} and
\isa{order{\isacharunderscore}bot} instances proved. Witness and
fixture facts are established constructively: concrete models are
exhibited as closed-form data and the substrate locales are
interpreted at them. The final theory,
\isa{Public{\isacharunderscore}Checker{\isacharunderscore}Witness},
exhibits concrete certificate/evidence pairs the verifier genuinely
accepts and fires all nine main theorems at them, so none of the
headline results is vacuously true.

\subsection*{Modelling notes}

The key type \isa{{\isacharprime}k} carries a \isa{linorder}
hypothesis in the run- and certificate-layer theorems: it enters
through the canonical clean-prefix construction, which enumerates
chunk-read domains with
\isa{sorted{\isacharunderscore}list{\isacharunderscore}of{\isacharunderscore}set}
and interleaves emitted events by source coordinate with a stable
\isa{sort{\isacharunderscore}key}, making the construction
deterministic. The four state-level continuation/restriction results
are constraint-free. Database primary keys are totally ordered in
practice, so the hypothesis does not narrow the intended
interpretation.

Source coordinates abstract positions in the source's totally ordered
change history (in deployments: a commit timestamp, log sequence
number, or binlog position). Chunk identifiers are plain naturals;
the run and certificate carriers are single-constructor datatypes
with named selectors.

\subsection*{Layer-to-theory map}

\begin{center}\small
\begin{tabular}{@{}l p{0.58\textwidth}@{}}
\toprule
Layer & Theories \\
\midrule
0 (source/replay substrate) & \isa{Source{\isacharunderscore}History}, \isa{Replay}, \isa{Scope}, \isa{Source{\isacharunderscore}Coordinates}, \isa{Virtual{\isacharunderscore}Cut{\isacharunderscore}State} --- the shared \isa{Dual{\isacharunderscore}Write{\isacharunderscore}Layer{\isadigit{0}}} parent session \\
0/1 (carrier-independent run core) & \isa{DBLog{\isacharunderscore}Run{\isacharunderscore}Core} \\
2--4 (carrier-independent cut core) & \isa{Virtual{\isacharunderscore}Cut{\isacharunderscore}Core} \\
1 (runs and wellformedness) & \isa{DBLog{\isacharunderscore}Run}, \isa{DBLog{\isacharunderscore}Run{\isacharunderscore}Substrate} \\
2 (run-side soundness) & \isa{DBLog{\isacharunderscore}Run{\isacharunderscore}Substrate{\isacharunderscore}Layer{\isadigit{2}}} \\
2--3 (virtual cuts, certificates) & \isa{Virtual{\isacharunderscore}Cut} \\
3--4 (certificate substrate locales) & \isa{DBLog{\isacharunderscore}Cert{\isacharunderscore}Substrate}(\isa{{\isacharunderscore}Inst}) \\
4 (whole table) & \isa{Layer{\isadigit{4}}{\isacharunderscore}Whole{\isacharunderscore}Table} \\
post-core extension & \isa{Continuation} \\
witnesses/fixtures & the \isa{{\isacharunderscore}Core}/\isa{{\isacharunderscore}Model}/\isa{{\isacharunderscore}Inst} families and the \isa{Layer{\isadigit{3}}{\isacharunderscore}Defect{\isacharunderscore}Regressions} regression theory \\
non-vacuity gate & \isa{Public{\isacharunderscore}Checker{\isacharunderscore}Witness} \\
\bottomrule
\end{tabular}
\end{center}

The witness/fixture families follow a uniform division of labor:
\isa{{\isacharunderscore}Core} theories hold carrier-independent data
(named coordinates, base states, source histories, expected values),
\isa{{\isacharunderscore}Model} theories exhibit small concrete
carrier realizations, and \isa{{\isacharunderscore}Inst} theories
interpret the substrate locales at those models and prove the
witness, boundary, and rejection facts.

\subsection*{Naming conventions}

\begin{center}\small
\begin{tabular}{ll}
\toprule
Convention & Meaning \\
\midrule
\isa{{\isacharunderscore}of}-suffixed accessors & run accessors (\isa{scope{\isacharunderscore}of}, \isa{frontier{\isacharunderscore}of}, \isa{clean{\isacharunderscore}prefix{\isacharunderscore}of}, \dots) \\
bare-named accessors & certificate accessors (\isa{scope}, \isa{frontier}, \isa{clean{\isacharunderscore}prefix}) \\
\isa{rr{\isacharunderscore}} / \isa{crr{\isacharunderscore}} prefixes & run-datatype / chunk-read-record field selectors \\
\isa{wit{\isacharunderscore}}, \isa{ex{\isacharunderscore}}, \isa{t{\isadigit{1}}{\isacharunderscore}}/\isa{t{\isadigit{2}}{\isacharunderscore}} prefixes & witness-family constants (also as suffixes: \isa{b{\isadigit{0}}{\isacharunderscore}w}, \isa{H{\isacharunderscore}ex}) \\
\isa{fix{\isacharunderscore}} prefix & negative/boundary fixture facts \\
\isa{layer{\isadigit{3}}{\isacharunderscore}} / \isa{l{\isadigit{4}}{\isacharunderscore}} prefixes & Layer-3 / Layer-4 fixture data and bundle facts \\
\isa{cf{\isacharunderscore}} prefix & certificate-substrate fixture-interpretation constants and facts \\
WF0--WF7 & the wellformedness clauses (WF3 is a deliberate gap) \\
\bottomrule
\end{tabular}
\end{center}

% generated text of all theories
\input{session}

% optional bibliography
\bibliographystyle{abbrv}
\bibliography{root}

\end{document}

%%% Local Variables:
%%% mode: latex
%%% TeX-master: t
%%% End:
